\origin{ai}
\subsection{Forced six-window Fibonacci cube closed-neighborhood Smith form}
\label{sec:window6-fibonacci-cube-closed-neighborhood-smith}

\paragraph{Finite carrier.}
Let $\{0,1\}^6$ be the set of all binary words $w=(w_1,w_2,w_3,w_4,w_5,w_6)$ of length six. Define
$$
\begin{aligned}
X_6=\{w\in \{0,1\}^6 : w_iw_{i+1}=0\hbox{ for every }1\leq i\leq 5\}.
\end{aligned}
$$
Thus $X_6$ is the finite set of length-six binary words with no adjacent $1$ symbols. The forced-window predicate is exactly the local exclusion of adjacent occupied positions; no limiting process, probability law, or outside model is used. The exact carrier size is
$$
\begin{aligned}
|X_6|=21.
\end{aligned}
$$
For bookkeeping one may partition $X_6$ by any fixed finite observable of the word, for example by its Hamming weight or by its final coordinate. Such partitions are only decompositions of the same twenty-one-word carrier; they do not add further hypotheses.

\paragraph{Finite relation.}
Define a graph $\Gamma_6$ on the vertex set $X_6$ by joining two distinct words when they differ in exactly one coordinate. Equivalently, $u,v\in X_6$ are adjacent precisely when
$$
\begin{aligned}
|\{i: u_i\neq v_i\}|=1.
\end{aligned}
$$
This relation is symmetric, irreflexive, and finite because it is a relation on the finite set $X_6$. The exact edge count for this graph is
$$
\begin{aligned}
|E(\Gamma_6)|=38,
\end{aligned}
$$
where edges are counted as undirected unordered pairs. This gives the full finite combinatorial object: twenty-one admissible words and thirty-eight Hamming-distance-one adjacencies among them.

\paragraph{Closed-neighborhood matrix.}
Choose any ordering of the twenty-one words of $X_6$. Let $A_6$ be the $21\times 21$ adjacency matrix of $\Gamma_6$ in that ordering: the entry in row $u$ and column $v$ is $1$ when $u$ and $v$ are adjacent, and is $0$ otherwise. Let $I_{21}$ be the $21\times 21$ identity matrix and put
$$
\begin{aligned}
N_6=I_{21}+A_6.
\end{aligned}
$$
Thus $N_6$ records closed neighborhoods: the diagonal $1$ records the vertex itself, while the off-diagonal $1$ entries record Hamming-distance-one neighbors that remain inside the no-adjacent-ones carrier. Reordering $X_6$ simultaneously permutes rows and columns, so the determinant up to the same value and the Smith invariant factors are intrinsic to the finite graph rather than to the chosen listing.

\paragraph{Exact arithmetic statement.}
The closed-neighborhood matrix has determinant
$$
\begin{aligned}
\det(N_6)=-144.
\end{aligned}
$$
Consequently $N_6$ has no rational kernel, and its image has finite index $144$ in $\mathbf{Z}^{21}$. Its Smith normal form has invariant factors
$$
\begin{aligned}
1,1,\ldots,1,144,
\end{aligned}
$$
with the factor $1$ occurring twenty times. Therefore the integral cokernel is the cyclic group
$$
\begin{aligned}
{\rm coker}(N_6:\mathbf{Z}^{21}\to \mathbf{Z}^{21})
  =\mathbf{Z}^{21}/N_6\mathbf{Z}^{21}\cong \mathbf{Z}/144\mathbf{Z}.
\end{aligned}
$$
This is the finite algebraic content of the construction: the forced six-window graph produces a closed-neighborhood operator whose only nontrivial Smith obstruction is a single cyclic factor of order $144$.

\paragraph{Endpoint-deleted witness extraction.}
There is also a distinguished unimodular minor. If the row and column indexed by the word $000001$ are deleted from $N_6$, the resulting $20\times 20$ integer matrix has determinant
$$
\begin{aligned}
1.
\end{aligned}
$$
This gives a concrete finite witness that the Smith form has at least twenty unit invariant factors: the endpoint-deleted block already contains a determinant-one minor of size twenty. Together with the determinant value $-144$ for the full matrix, this fixes the remaining invariant factor as $144$ and yields the cyclic cokernel displayed above.

\paragraph{Fibonacci alignment.}
The number $144$ also equals the twelfth Fibonacci number when the Fibonacci sequence is normalized by $F_0=0$, $F_1=1$, and $F_{n+2}=F_{n+1}+F_n$:
$$
\begin{aligned}
144=F_{12}.
\end{aligned}
$$
In this chapter the equality is recorded only as an exact integer alignment between the closed-neighborhood Smith obstruction and the Fibonacci sequence. The construction itself is fully finite: it starts from the twenty-one admissible binary words of length six, uses the Hamming-distance-one relation inside that set, and computes the integral invariants of $I_{21}+A_6$.

\paragraph{Claim boundary.}
No physical constant is identified by this computation. In particular, the chapter does not use or infer the fine-structure constant, does not treat $137$ or any physical value as an input or target, and does not assert numerical proximity to a measured quantity. No global flow law, external dependency, metrological comparison, or convention-dependent interpretation follows from the finite graph, its determinant, its Smith form, or the endpoint-deleted determinant-one block. Any such interpretation would require a separate certificate outside the finite construction stated here.
